% Generated draft: requires mathematical and source review.
\documentclass[11pt]{article}
\usepackage[margin=25mm]{geometry}
\usepackage{fontspec}
\setmainfont{Latin Modern Roman}
\usepackage{amsmath,amssymb,mathtools,unicode-math}
\setmathfont{Latin Modern Math}
\usepackage{enumitem,needspace,xurl,hyperref}
\setlength{\parindent}{0pt}
\setlength{\parskip}{5pt}
\setlength{\emergencystretch}{3em}
\newcommand{\N}{\mathbb{N}}
\newcommand{\Z}{\mathbb{Z}}
\newcommand{\Q}{\mathbb{Q}}
\newcommand{\R}{\mathbb{R}}
\newcommand{\C}{\mathbb{C}}
\newcommand{\F}{\mathbb{F}}
\newcommand{\A}{\mathbb{A}}
\newcommand{\PP}{\mathbb{P}}
\newcommand{\E}{\mathbb{E}}
\newcommand{\eps}{\varepsilon}
\newcommand{\dd}{\,\mathrm{d}}
\begin{document}
\section*{UnsolvedMath ID: 30001598}
\textbf{Source problem code:} catalogue entry
\subsection{Definitions and mathematical statement}
Let \(\pi\colon X\to\mathbb P^N_K\) denote the blowup at the finitely many points \(p_1,\ldots,p_r\) in the supplied statement. The ground field \(K\), the permitted dimensions, and the hypotheses on these points are not specified. Write \(H\) for the pullback of the hyperplane class and \(E_i\) for the exceptional divisors associated with the points, whenever this description applies. The divisor under consideration is
\[
D=dH-\sum_{i=1}^{r}m_iE_i.
\]
Interpreting \(D\) as an integral Cartier divisor, set
\[
h^0(X,kD)=\dim_K H^0\!\left(X,\mathcal O_X(kD)\right).
\]
The divisor \(D\) is called \emph{semi-effective} here if
\[
\exists k\in\mathbb Z_{>0}\qquad h^0(X,kD)>0.
\]
The question is whether there is a decision procedure which, given the point configuration and \(D\), terminates and determines whether this condition holds.

The source presents the following as an equivalent formulation: decide whether
\[
\exists t\in\mathbb Z_{>0}\ \exists\,0\ne F\in K[x_0,\ldots,x_N]_{td}
\quad\text{such that}\quad
\operatorname{ord}_{p_i}(F)\ge tm_i
\quad(1\le i\le r).
\]
Here \(K[x_0,\ldots,x_N]_{td}\) denotes the homogeneous forms of degree \(td\), when that degree is a nonnegative integer. Vanishing to order at least a nonnegative integer \(a\) means that a local representative of the section defined by \(F\) belongs to \(\mathfrak m_{p_i}^{a}\), where \(\mathfrak m_{p_i}\) is the maximal ideal at \(p_i\).

These formulations retain unresolved source qualifications: the integrality and sign restrictions on \(d,m_i\), the nature of the points, and the effective encoding of the field and geometric input are unstated. In particular, a precise algorithmic assertion requires an input class and an effective presentation; the record does not determine whether the field, dimension, or point configuration is fixed or variable. The claimed equivalence is retained as a source assertion and requires review under the intended hypotheses.
\subsection{Short English statement}
Given a divisor on projective space blown up at finitely many points, determine whether some positive integral multiple has a nonzero global section. The source also describes this as finding a nonzero homogeneous polynomial with proportionally scaled degree and prescribed vanishing orders. The field, admissible inputs, and several geometric hypotheses remain unspecified.
\subsection{Sources}
\begin{itemize}
\item Supplied local UnsolvedMath registry record, permanent ID \(30001598\), source problem code \texttt{catalogue entry}, titled \emph{Decidability of Semieffective Divisors on Blown-Up Projective Space}, in the Oberwolfach Reports Open Problems collection. No individual proposer is identified; the supplied proposed year is 2010.
\item Supplied report citation: \emph{Mini-Workshop: Linear Series on Algebraic Varieties (2011)}, Oberwolfach Reports, pp.\ 2613--2650, DOI \texttt{10.4171/owr/2010/45}. Supplied source URL and literature-source link, labelled \texttt{doi.org/10.4171/owr/2010/45}: \url{https://doi.org/10.4171/owr/2010/45}. The record attributes the problem to workshop open-problem discussions; it supplies no precise problem page or named contributor.
\item Literature assessment dated \texttt{2026-08-21}: the supplied curation labels the question open and reports that no general effective procedure in the stated setting was located. This is supplied dated curation, not an independently verified literature assessment. The record also labels the statement \texttt{corrected\_verified} and describes a manual semantic, typographic, and LaTeX review; these are supplied editorial claims, not independent mathematical verification.
\end{itemize}
\end{document}
